%% ──────────────────────────────────────────────────────────────
%% templates_latex/resume.tex — Jinja2 template that renders to LaTeX
%% ──────────────────────────────────────────────────────────────
%%
%% CUSTOM JINJA2 DELIMITERS
%% ────────────────────────
%% This template does NOT use standard {{ }}/{% %} Jinja2 syntax.
%% Instead it uses delimiters that look like LaTeX commands:
%%
%%          — insert a variable's value
%%           — end a logic block
%%          — end loop
%%          — template comment (stripped from output)
%%
%% WHY?  LaTeX uses { } everywhere (\textbf{bold}, \begin{itemize}).
%% If we used Jinja2's default {{ }}, the template engine would
%% misparse LaTeX braces as template tags, causing cryptic errors.
%% The \COMMAND{...} convention is unambiguous because no real LaTeX
%% command is named \VAR, \BLOCK, or \COMMENT.
%%
%% ALL USER DATA IS PRE-ESCAPED
%% ────────────────────────────
%% Before this template is rendered, every string from the user's
%% profile has already been run through escape_latex(), which
%% replaces LaTeX special characters (& % $ # _ { } ~ ^ \) with
%% safe equivalents (\& \% \$ etc.).  This means we can safely
%% inject  without worrying about compilation errors
%% from characters like & or #.
%% ──────────────────────────────────────────────────────────────

\documentclass[11pt, letterpaper]{article}

%% ── Packages ─────────────────────────────────────────────────
\usepackage[utf8]{inputenc}          % UTF-8 input encoding
\usepackage[T1]{fontenc}             % proper output glyph encoding
\usepackage[scaled=0.92]{helvet}     % Helvetica font (clean sans-serif)
\renewcommand{\familydefault}{\sfdefault} % Set sans-serif as default
\usepackage[margin=0.5in]{geometry}  % tight but readable margins
\usepackage{enumitem}                % customise list spacing
\usepackage{titlesec}                % customise section headings
\usepackage{hyperref}                % clickable links
\usepackage{xcolor}                  % colours for links
\usepackage{marvosym}            % Classic icons for contact info
\usepackage[normalem]{ulem}          % For clean underlines

%% ── Hyperlink styling ────────────────────────────────────────
\hypersetup{
    colorlinks=true,
    linkcolor=black,
    urlcolor=blue!70!black,
    pdfborder={0 0 0},
}

%% ── Section heading format ───────────────────────────────────
%% \titleformat{command}{format}{label}{sep}{before-code}[after-code]
%% We use a simple bold heading (without small caps) with a rule underneath.
\titleformat{\section}
  {\large\bfseries\raggedright}          % format
  {}                                      % label (none)
  {0em}                                   % separator
  {}                                      % before-code
  [\titlerule]                            % after-code: horizontal rule

\titlespacing*{\section}{0pt}{10pt}{4pt}

%% ── Remove page numbers ─────────────────────────────────────
\pagestyle{empty}

%% ── Tighter list spacing ────────────────────────────────────
\setlist[itemize]{label=\textbullet, nosep, leftmargin=1.5em, topsep=2pt}

%% ── Zero Paragraph Indentation ──────────────────────────────
\setlength{\parindent}{0pt}

%% ── Unicode Rupee Symbol Mapping ────────────────────────────
\usepackage{tfrupee}
\usepackage{newunicodechar}
\newunicodechar{₹}{\rupee}

%% ══════════════════════════════════════════════════════════════
%%  DOCUMENT BEGIN
%% ══════════════════════════════════════════════════════════════
\begin{document}

%% ── Header ───────────────────────────────────────────────────
\begin{center}
    {\LARGE\bfseries John Developer}\\[5pt]
    \small
    %
        \Mobilefone\ 123-456-7890%
    %
    %
         \quad %
        \Letter\ \href{mailto:john.dev@example.com}{john.dev@example.com}%
    %
    %
        \quad \textbullet\ \href{https://github.com/johndev}{GitHub}%
    %
        \quad \textbullet\ \href{https://example.com}{Portfolio}%
    
\end{center}

\vspace{-4pt}


%
%% ── Education ────────────────────────────────────────────────
\section{Education}
%
\textbf{Kalinga Institute of Industrial College} \hfill Aug 2024 -- May 2028 \\
\textit{B.Tech} \textit{(CGPA: 8.38)}%
%
 \hfill CSE (AI-ML)%
%
%
%



%
%% ── Experience ───────────────────────────────────────────────
\section{Experience}
%
\textbf{Swastik Agency} \hfill Jan 2024 -- Present \\
\textit{HouseHusband} \ | \ \textit{Python}%
%
 \hfill onsite%
%

\begin{itemize}[label=\textasteriskcentered, leftmargin=1.5em]

    \item made best cookies  in the world.

    \item made icereams.

    \item learnt gardening and became vegetarian.

\end{itemize}
%

\vspace{4pt}
%
%
\textbf{Tech Corp} \hfill 2022 -- Present \\
\textit{Junior Engineer} \ | \ \textit{Artificial Intelligence (AI), Software Infrastructure}%
%
 \hfill Hybrid%
%

\begin{itemize}[label=\textasteriskcentered, leftmargin=1.5em]

    \item Developed and maintained web applications utilizing React, contributing to front-end development initiatives.

\end{itemize}
%
%
%



%
%% ── Projects ─────────────────────────────────────────────────
\section{Projects}
%
\textbf{Competitive Programming Ally} \ | \ \href{https://github.com/0xPolybit/cp-ally-ide}{\textbf{(GitHub)}} \ | \ \textit{Java, Java Swing, File I/O}%
%
\\ you can do codeforeces here it self%
%

\begin{itemize}[label=\textasteriskcentered, leftmargin=1.5em]

    \item Developed a specialized IDE for competitive programming that automates the process of fetching and testing sample test cases.

    \item Implemented a custom testing engine for real-time output comparison, significantly reducing debugging time during contests.

    \item Designed a streamlined user interface using Java Swing to minimize context switching and improve coding efficiency.

\end{itemize}
%
%
%



%
%% ── Skills ───────────────────────────────────────────────────
\section{Technical Skills}
%
\textbf{Artificial Intelligence:} LLMsLLMs \\ %
%
\textbf{Programming Languages:} JavaScript, TypeScript, JAVA, Python \\ %
%
\textbf{Software and Tools:} Git, GitHub \\ %
%
\textbf{Web Frameworks:} React.js, Next.js, Vite.js%
%






%
%% ── Achievements ─────────────────────────────────────────────
\section{Achievements}
Secured 2nd place in the SDIS Competitive Programming Contest, winning a cash prize of ₹5,000


\end{document}
